BARTH, Adam; JACKSON, Collin; MITCHELL, John C. Robust defenses for cross-site request forgery. In: ACM CONFERENCE ON COMPUTER AND COMMUNICATIONS SECURITY (CCS), 15., 2008, Alexandria. Proceedings... New York: ACM, 2008. p. 75–88. DOI: 10.1145/1455770.1455782. Disponível em: https://dl.acm.org/doi/10.1145/1455770.1455782.
DOMÍNGUEZ-GARCÍA, Antonio De Jesús; LIMÓN, Xavier; OCHARÁN-HERNÁNDEZ, Jorge Octavio; PÉREZ-ARRIAGA, Juan Carlos. Security testing for web applications: a systematic literature review. In: IEEE CONFERENCE, 2024. Proceedings... [S. l.]: IEEE, 2024. Disponível em: https://ieeexplore.ieee.org/document/10568290. Acesso em: fev. 2026.
DOUPE, Adam; COVA, Marco; VIGNA, Giovanni. Why Johnny can't pentest: an analysis of black-box web vulnerability scanners. In: INTERNATIONAL CONFERENCE ON DETECTION OF INTRUSIONS AND MALWARE, AND VULNERABILITY ASSESSMENT (DIMVA), 7., 2010, Bonn. Proceedings... [S. l.]: Springer, 2010. p. 111–131. (Lecture Notes in Computer Science, v. 6201). DOI: 10.1007/978-3-642-14215-4_7. Disponível em: <https://link.springer.com/chapter/10.1007/978-3-642-14215-4_7>.
GUPTA, Shashank; GUPTA, B. B. Cross-site scripting (XSS) attacks and defense mechanisms: classification and state-of-the-art. International Journal of System Assurance Engineering and Management, [S. l.]: Springer, v. 8, p. 512–530, 2015. DOI: 10.1007/s13198-015-0376-0. Disponível em: https://link.springer.com/article/10.1007/s13198-015-0376-0.
HALFOND, William G. J.; VIEGAS, Jeremy; ORSO, Alessandro. A classification of SQL-injection attacks and countermeasures. In: IEEE INTERNATIONAL SYMPOSIUM ON SECURE SOFTWARE ENGINEERING (ISSSE), 2006. Proceedings... [S. l.]: IEEE, 2006. Disponível em: <https://www.researchgate.net/publication/249773840_A_Classification_of_SQL_Injection_Attacks_and_Countermeasures>. Acesso em: fev. 2026.
HEIDERICH, Mario et al. mXSS attacks: attacking well-secured web-applications by using innerHTML mutations. In: ACM SIGSAC CONFERENCE ON COMPUTER AND COMMUNICATIONS SECURITY (CCS), 2013. Proceedings... New York: ACM, 2013. p. 777–788. DOI: 10.1145/2508859.2516723. Disponível em: https://dl.acm.org/doi/10.1145/2508859.2516723.
OPEN WEB APPLICATION SECURITY PROJECT (OWASP). A05:2021 – Security Misconfiguration. In: OWASP top 10:2021. [S. l.]: OWASP, 2021. Disponível em: <https://owasp.org/Top10/2021/A05_2021-Security_Misconfiguration/>. Acesso em: fev. 2026.
SACRAMENTO, Lucas S. C. et al. Automação de autenticação para testes de segurança em aplicações web no ZAP. In: SIMPÓSIO BRASILEIRO DE CIBERSEGURANÇA (SBSeg), 24., 2024, São José dos Campos/SP. Anais... Porto Alegre: SBC, 2024. p. 760-766. DOI: 10.5753/sbseg.2024.241760. Disponível em: https://sol.sbc.org.br/index.php/sbseg/article/view/30065.
SETH, Aishwarya et al. Comparing effectiveness and efficiency of interactive application security testing (IAST) and runtime application self-protection (RASP) tools in a large Java-based system. Empirical Software Engineering, [S. l.]: Springer, v. 30, artigo 67, 2025. DOI: 10.1007/s10664-025-10621-5. Disponível em: https://link.springer.com/article/10.1007/s10664-025-10621-5. Acesso em: fev. 2026.
SUREDA RIERA, Tomás et al. Systematic approach for web protection runtime tools' effectiveness analysis. Computer Modeling in Engineering & Sciences, [S. l.], v. 133, n. 3, 2022. DOI: 10.32604/cmes.2022.020976. Disponível em: https://www.techscience.com/CMES/v133n3/49215/html. Acesso em: fev. 2026.
